% Standalone research note for the hybrid-local-solver project.
%
% This note bounds the seed diagonal of the discounted PageRank Green kernel
% in terms of ordinary seed degree and graph volume, then applies the bound to
% the bounded-seed-degree question behind signed one-hop class lower bounds.

\documentclass[11pt]{article}

% Common class-neutral preamble for standalone notes under manuscript/notes/.
% Note entry points all live exactly two directories below manuscript/.

\usepackage[margin=1in]{geometry}
\usepackage{amsmath,amssymb,amsthm,mathtools,bm}
\usepackage{algorithm}
\usepackage{algorithmic}
\usepackage{booktabs,tabularx,multirow,array,longtable}
\usepackage{enumitem}
\usepackage{microtype}
\usepackage[numbers,sort&compress]{natbib}
\usepackage{xcolor}
\usepackage[colorlinks=true,citecolor=blue,linkcolor=black,urlcolor=blue]{hyperref}
\usepackage[capitalize,nameinlink,noabbrev]{cleveref}

\newtheorem{theorem}{Theorem}[section]
\newtheorem{lemma}[theorem]{Lemma}
\newtheorem{proposition}[theorem]{Proposition}
\newtheorem{corollary}[theorem]{Corollary}
\newtheorem{conjecture}[theorem]{Conjecture}
\theoremstyle{definition}
\newtheorem{definition}[theorem]{Definition}
\newtheorem{assumption}[theorem]{Assumption}
\newtheorem{problem}[theorem]{Problem}
\newtheorem{observation}[theorem]{Observation}
\theoremstyle{remark}
\newtheorem{remark}[theorem]{Remark}
\theoremstyle{plain}

% Canonical mathematical notation for the active manuscript.
%
% This file preserves the recurring notation style used in the author's
% NeurIPS 2024 and 2025 papers while excluding unrelated tensor and
% machine-learning boilerplate from those archived command collections.
%
% Prefix convention:
%   \v...  bold vectors
%   \m...  bold matrices
%   \g...  calligraphic sets, graphs, and families
%   \s...  blackboard-bold spaces and number systems

\newcommand{\method}{\textsc{HybridLocalSolver}}

% Font and scalar shorthands retained from the archived manuscripts.
\newcommand{\mc}{\mathcal}
\newcommand{\R}{\mathbb{R}}
\newcommand{\E}{\mathbb{E}}
\newcommand{\G}{\mathbb{G}}
\newcommand{\N}{\mathcal{N}}
% Accuracy namespaces are semantic: do not add a bare \eps alias.
% Degree-normalized residual tolerance of classical APPR is kept distinct from
% the objective-gap and proximal fixed-point tolerances of the RPPR sections.
\newcommand{\epsappr}{\varepsilon_{\mathrm{appr}}}
\newcommand{\epsobj}{\varepsilon_{\mathrm{obj}}}
\newcommand{\epspg}{\varepsilon_{\mathrm{pg}}}
\newcommand{\epsppr}{\varepsilon_{\mathrm{ppr}}}
\newcommand{\epsin}{\varepsilon_{\mathrm{in}}}
\newcommand{\epskkt}{\varepsilon_{\mathrm{kkt}}}
\newcommand{\epsburn}{\varepsilon_{\mathrm{burn}}}
\newcommand{\epssol}{\varepsilon_{\mathrm{sol}}}
\newcommand{\epspprinst}[1]{\varepsilon_{\mathrm{ppr},#1}}
\newcommand{\trans}{\mathsf{T}}
\newcommand{\grad}{\nabla}

% Delimiters and generic helpers.
% Norms and inner products use fixed-size delimiters by default.
% Use an optional size (e.g. \norm[\big]{...}) or * for automatic sizing.
\DeclarePairedDelimiter{\norm}{\lVert}{\rVert}
\newcommand{\abs}[1]{\left\lvert #1 \right\rvert}
\DeclarePairedDelimiterX{\inner}[2]{\langle}{\rangle}{#1, #2}
\newcommand{\ip}[2]{\inner{#1}{#2}}
\newcommand{\card}[1]{\left\lvert #1 \right\rvert}
\newcommand{\set}[1]{\left\{#1\right\}}
\newcommand{\ceil}[1]{\left\lceil #1 \right\rceil}
\newcommand{\floor}[1]{\left\lfloor #1 \right\rfloor}
\newcommand{\comp}[1]{\overline{#1}}
\newcommand{\conj}[1]{\overline{#1}}
\newcommand{\mean}[1]{\overline{#1}}
\newcommand{\bigO}{\mathcal{O}}
\newcommand{\softO}{\widetilde{\mathcal{O}}}
\newcommand{\softOmega}{\widetilde{\Omega}}
\newcommand{\pos}[1]{\left[#1\right]_{+}}

% Bold lowercase vectors.
\newcommand{\va}{\bm{a}}
\newcommand{\vb}{\bm{b}}
\newcommand{\vc}{\bm{c}}
\newcommand{\vd}{\bm{d}}
\newcommand{\ve}{\bm{e}}
\newcommand{\vf}{\bm{f}}
\newcommand{\vg}{\bm{g}}
\newcommand{\vh}{\bm{h}}
\newcommand{\vi}{\bm{i}}
\newcommand{\vj}{\bm{j}}
\newcommand{\vk}{\bm{k}}
\newcommand{\vl}{\bm{l}}
\newcommand{\vm}{\bm{m}}
\newcommand{\vn}{\bm{n}}
\newcommand{\vo}{\bm{o}}
\newcommand{\vp}{\bm{p}}
\newcommand{\vq}{\bm{q}}
\newcommand{\vr}{\bm{r}}
\newcommand{\vs}{\bm{s}}
\newcommand{\vt}{\bm{t}}
\newcommand{\vu}{\bm{u}}
\newcommand{\vv}{\bm{v}}
\newcommand{\vw}{\bm{w}}
\newcommand{\vx}{\bm{x}}
\newcommand{\vy}{\bm{y}}
\newcommand{\vz}{\bm{z}}

% Bold Greek vectors and special vectors.
\newcommand{\valpha}{\bm{\alpha}}
\newcommand{\vbeta}{\bm{\beta}}
\newcommand{\vdelta}{\bm{\delta}}
\newcommand{\vepsilon}{\bm{\epsilon}}
\newcommand{\veta}{\bm{\eta}}
\newcommand{\vgamma}{\bm{\gamma}}
\newcommand{\vlambda}{\bm{\lambda}}
\newcommand{\vmu}{\bm{\mu}}
\newcommand{\vomega}{\bm{\omega}}
\newcommand{\vphi}{\bm{\phi}}
\newcommand{\vpi}{\bm{\pi}}
\newcommand{\vpsi}{\bm{\psi}}
\newcommand{\vrho}{\bm{\rho}}
\newcommand{\vsigma}{\bm{\sigma}}
\newcommand{\vtheta}{\bm{\theta}}
\newcommand{\vxi}{\bm{\xi}}
\newcommand{\vell}{\bm{\ell}}
\newcommand{\vDelta}{\bm{\Delta}}
\newcommand{\vGamma}{\bm{\Gamma}}
\newcommand{\vLambda}{\bm{\Lambda}}
\newcommand{\vPhi}{\bm{\Phi}}
\newcommand{\vSigma}{\bm{\Sigma}}
\newcommand{\vzero}{\bm{0}}
\newcommand{\vone}{\bm{1}}
\newcommand{\one}{\mathbf{1}}
\newcommand{\zero}{\vzero}
\newcommand{\eunit}[1]{\bm{e}_{#1}}
\newcommand{\vDeltax}{\bm{\Delta}\bm{x}}

% Bold uppercase matrices.
\newcommand{\mA}{\bm{A}}
\newcommand{\mB}{\bm{B}}
\newcommand{\mC}{\bm{C}}
\newcommand{\mD}{\bm{D}}
\newcommand{\mE}{\bm{E}}
\newcommand{\mF}{\bm{F}}
\newcommand{\mG}{\bm{G}}
\newcommand{\mH}{\bm{H}}
\newcommand{\mI}{\bm{I}}
\newcommand{\mJ}{\bm{J}}
\newcommand{\mK}{\bm{K}}
\newcommand{\mL}{\bm{L}}
\newcommand{\mM}{\bm{M}}
\newcommand{\mN}{\bm{N}}
\newcommand{\mO}{\bm{O}}
\newcommand{\mP}{\bm{P}}
\newcommand{\mQ}{\bm{Q}}
\newcommand{\mR}{\bm{R}}
\newcommand{\mS}{\bm{S}}
\newcommand{\mT}{\bm{T}}
\newcommand{\mU}{\bm{U}}
\newcommand{\mV}{\bm{V}}
\newcommand{\mW}{\bm{W}}
\newcommand{\mX}{\bm{X}}
\newcommand{\mY}{\bm{Y}}
\newcommand{\mZ}{\bm{Z}}

% Bold Greek matrices retained from the graph papers.
\newcommand{\mBeta}{\bm{\beta}}
\newcommand{\mDelta}{\bm{\Delta}}
\newcommand{\mGamma}{\bm{\Gamma}}
\newcommand{\mLambda}{\bm{\Lambda}}
\newcommand{\mOmega}{\bm{\Omega}}
\newcommand{\mPhi}{\bm{\Phi}}
\newcommand{\mPi}{\bm{\Pi}}
\newcommand{\mSigma}{\bm{\Sigma}}
\newcommand{\mTheta}{\bm{\Theta}}

% Calligraphic symbols.
\newcommand{\gA}{\mathcal{A}}
\newcommand{\gB}{\mathcal{B}}
\newcommand{\gC}{\mathcal{C}}
\newcommand{\gD}{\mathcal{D}}
\newcommand{\gE}{\mathcal{E}}
\newcommand{\gF}{\mathcal{F}}
\newcommand{\gG}{\mathcal{G}}
\newcommand{\gH}{\mathcal{H}}
\newcommand{\gI}{\mathcal{I}}
\newcommand{\gJ}{\mathcal{J}}
\newcommand{\gK}{\mathcal{K}}
\newcommand{\gL}{\mathcal{L}}
\newcommand{\gM}{\mathcal{M}}
\newcommand{\gN}{\mathcal{N}}
\newcommand{\gO}{\mathcal{O}}
\newcommand{\gP}{\mathcal{P}}
\newcommand{\gQ}{\mathcal{Q}}
\newcommand{\gR}{\mathcal{R}}
\newcommand{\gS}{\mathcal{S}}
\newcommand{\gT}{\mathcal{T}}
\newcommand{\gU}{\mathcal{U}}
\newcommand{\gV}{\mathcal{V}}
\newcommand{\gW}{\mathcal{W}}
\newcommand{\gX}{\mathcal{X}}
\newcommand{\gY}{\mathcal{Y}}
\newcommand{\gZ}{\mathcal{Z}}

% Blackboard-bold symbols.
\newcommand{\sA}{\mathbb{A}}
\newcommand{\sB}{\mathbb{B}}
\newcommand{\sC}{\mathbb{C}}
\newcommand{\sD}{\mathbb{D}}
\newcommand{\sE}{\mathbb{E}}
\newcommand{\sF}{\mathbb{F}}
\newcommand{\sG}{\mathbb{G}}
\newcommand{\sH}{\mathbb{H}}
\newcommand{\sI}{\mathbb{I}}
\newcommand{\sJ}{\mathbb{J}}
\newcommand{\sK}{\mathbb{K}}
\newcommand{\sL}{\mathbb{L}}
\newcommand{\sM}{\mathbb{M}}
\newcommand{\sN}{\mathbb{N}}
\newcommand{\sO}{\mathbb{O}}
\newcommand{\sP}{\mathbb{P}}
\newcommand{\sQ}{\mathbb{Q}}
\newcommand{\sR}{\mathbb{R}}
\newcommand{\sS}{\mathbb{S}}
\newcommand{\sT}{\mathbb{T}}
\newcommand{\sU}{\mathbb{U}}
\newcommand{\sV}{\mathbb{V}}
\newcommand{\sW}{\mathbb{W}}
\newcommand{\sX}{\mathbb{X}}
\newcommand{\sY}{\mathbb{Y}}
\newcommand{\sZ}{\mathbb{Z}}

% Frequently used operators.
\DeclareMathOperator{\Cov}{Cov}
\DeclareMathOperator{\Var}{Var}
\DeclareMathOperator{\diag}{diag}
\DeclareMathOperator{\nei}{Nei}
\DeclareMathOperator{\nnz}{nnz}
\DeclareMathOperator{\pr}{pr}
\DeclareMathOperator{\prox}{prox}
\DeclareMathOperator{\push}{push}
\DeclareMathOperator{\res}{res}
\DeclareMathOperator{\rel}{Rel}
\DeclareMathOperator{\relu}{ReLU}
\DeclareMathOperator{\sign}{sign}
\DeclareMathOperator{\supp}{supp}
\DeclareMathOperator{\tr}{tr}
\DeclareMathOperator{\Tr}{Tr}
\DeclareMathOperator{\vol}{vol}
\DeclareMathOperator{\work}{work}
\DeclareMathOperator{\Work}{Work}
\DeclareMathOperator*{\argmax}{arg\,max}
\DeclareMathOperator*{\argmin}{arg\,min}

% Project-wide names and claim-status labels used by standalone research notes.
% Mathematical notation belongs in math_commands.tex.

% Algorithms and solver families.
\newcommand{\AESP}{\textsc{AESP}}
\newcommand{\APPR}{\textsc{APPR}}
\newcommand{\ASPR}{\textsc{ASPR}}
\newcommand{\FISTA}{\textsc{FISTA}}
\newcommand{\ISTA}{\textsc{ISTA}}
\newcommand{\LOCAPPR}{\textsc{LocAPPR}}
\newcommand{\LOCGD}{\textsc{LocGD}}
\newcommand{\LOCSOR}{\textsc{LocSOR}}
\newcommand{\LOCCH}{\textsc{LocCH}}
\newcommand{\LOCHB}{\textsc{LocHB}}
\newcommand{\LocProxGS}{\textsc{LocProxGS}}
\newcommand{\Hybrid}{\textsc{AESP--LocSOR}}

% Claim classifications required by the repository research-note protocol.
\newcommand{\statussource}{\textbf{Source result}}
\newcommand{\statusproved}{\textbf{Proved here}}
\newcommand{\statusconditional}{\textbf{Conditional}}
\newcommand{\statusconjecture}{\textbf{Open / conjectural}}
\newcommand{\statusopen}{\textbf{Open}}
\newcommand{\statusempirical}{\textbf{Empirical}}
\newcommand{\statusobservation}{\textbf{Empirical observation}}
\newcommand{\statusrefuted}{\textbf{Refuted route}}

% Compatibility names used by the older complete-note presentation layer.
\newcommand{\Status}[2]{\textbf{#1:} #2}
\newcommand{\SourceStatus}{\textnormal{\textsc{Source result}}}
\newcommand{\ProvedStatus}{\textnormal{\textsc{Proved here}}}
\newcommand{\ConditionalStatus}{\textnormal{\textsc{Conditional}}}
\newcommand{\ComputedStatus}{\textnormal{\textsc{Computed}}}
\newcommand{\RefutedStatus}{\textnormal{\textsc{Refuted route}}}
\newcommand{\OpenStatus}{\textnormal{\textsc{Open}}}



\title{Bounded Seed Degree Forbids Star-Scale\
Self-Return Amplification}
\author{Baojian Zhou\\
\small Working research note for \texttt{hybrid-local-solver}}
\date{August 2026}

\begin{document}
\maketitle

\begin{abstract}
For source-aligned personalized PageRank seeded at $v$, put
$c_{\alpha}=(1-\alpha)/(1+\alpha)$ and
$\gamma_{\alpha}=1-c_{\alpha}$.  We prove the graph-uniform bound
\[
 \frac{\pi_v}{\gamma_{\alpha}}
 \le
 \frac{d_v}{\vol(G)\gamma_{\alpha}}
 +\frac{1}{1+c_{\alpha}}
 +\frac{2\pi d_v}{\sqrt{c_{\alpha}\gamma_{\alpha}}}.
\]
The proof is a local normalized-Laplacian spectral estimate: the mass at $v$
of nonstationary eigenvectors below frequency $t$ is at most
$4d_v\sqrt t$.  A Stieltjes integral then separates finite-graph stationarity
from the universal one-dimensional return envelope.

If the semantic support
$S_{\epsppr}=\{u:\pi_u/d_u>\epsppr\}$ has volume
$\Omega(1/\epsppr)$, the bound implies
$\alpha\pi_v/\gamma_{\alpha}\to0$ along every joint limit
$\alpha,\epsppr\to0$ with bounded ordinary seed degree.  Thus no such family
has the star's $\Theta(1/\alpha)$ seed self-return amplification.  More
quantitatively, retaining that amplification requires
$d_v=\Omega(1/(\epsppr+\sqrt\alpha))$, and the stronger
$\Omega(1/\epsppr)$ conclusion needs the regime
$\alpha=O(\epsppr^2)$.  This closes only the self-return mechanism: it neither
rules out a distributed signed-relaxation lower bound nor proves an
algorithmic upper bound.
\end{abstract}

% Canonical source-aligned PageRank/RPPR reference formulation.
%
% This fragment is intentionally heading-free at the section level so that the
% active paper and every standalone note can input it. It is a manuscript
% reference convention, not an implementation-wide residual decision.
% Primary source: Fountoulakis and Martinez-Rubio, arXiv:2602.21138v2,
% PDF p. 3, Section 3 and equation (RPPR).

\subsection{Common source-aligned PageRank and RPPR model}
\label{subsec:shared-source-aligned-problem}

All documents in this manuscript workspace use the following reference
language. It follows the plain italic notation of the comparison paper:
\(x,s,A,D,Q,I\) denote column vectors or matrices as determined by their
displayed domains. This is the scoped exception to the project's usual
bold-vector and bold-matrix typography. A note may impose a stronger
assumption after this block, but it must not redefine these objects.

Let \(G=(V,E)\) be a finite, simple, undirected, unweighted graph without
isolated vertices, let \(V=[n]:=\{1,\ldots,n\}\), and let
\(A\in\{0,1\}^{n\times n}\) be its symmetric adjacency matrix. For
\(i\in V\), write
\[
    \mathcal N(i):=\{j\in V:\{i,j\}\in E\},
    \qquad
    d_i:=|\mathcal N(i)|,
    \qquad
    D:=\diag(d_1,\ldots,d_n).
\]
For \(S\subseteq V\), define
\[
    \mathcal N(S):=\bigcup_{i\in S}\mathcal N(i),
    \qquad
    \partial S:=\mathcal N(S)\setminus S,
    \qquad
    \operatorname{Ext}(S):=V\setminus(S\cup\partial S),
\]
\begin{equation}
    \vol(S):=\sum_{i\in S}d_i.
    \label{eq:shared-volume}
\end{equation}
The support of a vector is
\(\supp(x):=\{i\in[n]:x_i\neq0\}\). Matrix inequalities use the Loewner
order, and all unqualified vector inequalities are coordinatewise.

The seed is a column distribution
\begin{equation}
    s\in\R_+^n,
    \qquad
    \inner{\one}{s}=1.
    \label{eq:shared-seed}
\end{equation}
The single-seed specialization is always written \(s=e_v\), where \(v\in V\)
is the seed vertex. Thus \(s\) is never used as a vertex label. For
\(\alpha\in(0,1]\), define
\begin{equation}
\begin{aligned}
    \mathcal L&:=I-D^{-1/2}AD^{-1/2},\\
    Q&:=\alpha I+\frac{1-\alpha}{2}\mathcal L
      =\frac{1+\alpha}{2}I
       -\frac{1-\alpha}{2}D^{-1/2}AD^{-1/2},\\
    b&:=\alpha D^{-1/2}s.
\end{aligned}
    \label{eq:shared-pagerank-matrices}
\end{equation}
Since \(0\preceq\mathcal L\preceq2I\),
\(\alpha I\preceq Q\preceq I\). The shared smooth PageRank quadratic is
\begin{equation}
    f(x):=\frac12\inner{x}{Qx}-\inner{b}{x},
    \qquad
    \nabla f(x)=Qx-b.
    \label{eq:shared-pagerank-objective}
\end{equation}
Its unique unregularized minimizer and associated lazy personalized PageRank
vector are
\begin{equation}
    x^0:=Q^{-1}b,
    \qquad
    \pi:=D^{1/2}x^0.
    \label{eq:shared-ppr-solution}
\end{equation}

For a regularization parameter \(\rho>0\), reserve \(g_\rho\) for the
nonsmooth weighted \(\ell_1\) term and define
\begin{equation}
    g_\rho(x):=\alpha\rho\norm{D^{1/2}x}_1,
    \qquad
    F_\rho(x):=f(x)+g_\rho(x),
    \qquad
    x^\star(\rho):=\argmin_{x\in\R^n}F_\rho(x).
    \label{eq:shared-rppr-objective}
\end{equation}
When \(\rho\) is fixed, \(x^\star\) abbreviates \(x^\star(\rho)\), never the
unregularized point \(x^0\). Write
\[
    S^\star(\rho):=\supp(x^\star(\rho)),
    \qquad
    I^\star(\rho):=[n]\setminus S^\star(\rho).
\]
The source records \(x^\star(\rho)\geq0\),
\(\vol(S^\star(\rho))\leq1/\rho\), and the coordinatewise conditions
\begin{equation}
    \nabla_i f(x^\star(\rho))
    \in
    \begin{cases}
        \{-\alpha\rho\sqrt{d_i}\}, & x_i^\star(\rho)>0,\\
        [-\alpha\rho\sqrt{d_i},0], & x_i^\star(\rho)=0.
    \end{cases}
    \label{eq:shared-rppr-kkt}
\end{equation}

The common computational unit is one repeated adjacency-list scan: scanning
coordinate \(i\) costs \(d_i\), and scanning a set \(S\) costs \(\vol(S)\).
Iteration count and wall-clock time do not replace this degree-weighted work
measure.

Finally, accuracy symbols remain distinct. The APPR threshold is
\(\epsappr\); \(\epsobj\) denotes an objective-gap target;
\(\epspg\) denotes the source experiment's proximal fixed-point tolerance; and
\(\epsppr\) may denote a document-scoped degree-normalized PPR error target.
No equality or conversion among these quantities is assumed. In particular,
this shared model does not resolve the repository-wide residual decision.


\section{Scope and the diagonal return quantity}
\label{sec:bounded-return-scope}

\paragraph{Claim status.}
All lemmas, theorems, and corollaries are \statusproved{} in the shared finite,
simple, connected, undirected, unweighted exact-arithmetic model.  The family
screen in \cref{sec:bounded-return-families} combines proved formulas with a
separately labeled numerical check.  No algorithmic lower bound is inferred
from a large or small Green diagonal.

Specialize the shared PageRank definition to $s=e_v$ and write
\begin{equation}
 c:=c_{\alpha}:=\frac{1-\alpha}{1+\alpha},
 \qquad
 \gamma:=\gamma_{\alpha}:=1-c=\frac{2\alpha}{1+\alpha},
 \qquad
 H:=I-cAD^{-1}.
 \label{eq:bounded-return-parameters}
\end{equation}
Then $H\pi=\gamma e_v$.  The scalar screened in this note is
\begin{equation}
 \mathcal A_v:=\frac{\pi_v}{\gamma}=(H^{-1})_{vv}.
 \label{eq:bounded-return-amplification}
\end{equation}
It is the geometrically discounted expected number of visits to $v$ by the
ordinary random walk started at $v$.  The desired star scaling is
$\mathcal A_v=\Theta(1/\alpha)$.  Throughout, $d_v$ means the ordinary
combinatorial degree, equivalently the adjacency-list scan cost.  It is not
normalized by the minimum or average degree.

Define the semantic level set
\begin{equation}
 S_{\epsppr}:=\left\{u\in V:\frac{\pi_u}{d_u}>\epsppr\right\}.
 \label{eq:bounded-return-semantic-support}
\end{equation}
The asymptotic question is joint: both $\alpha\to0$ and $\epsppr\to0$ while
$\vol(S_{\epsppr})=\Theta(1/\epsppr)$.  At fixed accuracy on a fixed finite
graph, the stationary mode can yield $\mathcal A_v=\Theta(1/\alpha)$ even at
bounded degree; that different limit is not ruled out.

\section{A local low-frequency spectral bound}
\label{sec:bounded-return-spectral}

Let the shared normalized Laplacian $L$ have orthonormal eigenpairs
$(\lambda_k,\varphi_k)$ ordered as
\[
 0=\lambda_1<\lambda_2\le\cdots\le\lambda_n\le2.
\]
Connectedness gives
\begin{equation}
 \varphi_1=\frac{D^{1/2}\one}{\sqrt{\vol(G)}},
 \qquad
 \varphi_1(v)^2=\frac{d_v}{\vol(G)}.
 \label{eq:bounded-return-stationary-mode}
\end{equation}
For $t>0$, define the nonstationary local spectral mass
\begin{equation}
 F_v(t):=\sum_{\substack{k\ge2\\\lambda_k\le t}}\varphi_k(v)^2.
 \label{eq:bounded-return-spectral-mass}
\end{equation}

\begin{lemma}[Local spectral-measure envelope; \statusproved]
\label{lem:bounded-return-spectral-measure}
For every $v\in V$ and every $t\in(0,2]$,
\begin{equation}
 F_v(t)\le4d_v\sqrt t.
 \label{eq:bounded-return-spectral-envelope}
\end{equation}
\end{lemma}

\begin{proof}
It suffices to consider $t\le1$, since for $t>1$ one has
$F_v(t)\le1\le4d_v\sqrt t$.  Let $P_t$ be the orthogonal projector onto the
span of eigenvectors with eigenvalues in $(0,t]$, and put
\[
 f:=P_te_v,
 \qquad F:=\|f\|_2^2=F_v(t),
 \qquad g:=D^{-1/2}f.
\]
If $F=0$ there is nothing to prove.  Since $P_t$ is a projector,
\begin{equation}
 f_v=\langle e_v,P_te_v\rangle=\|P_te_v\|_2^2=F,
 \qquad
 g_v=\frac{F}{\sqrt{d_v}}.
 \label{eq:bounded-return-projector-at-v}
\end{equation}
The projector excludes the stationary eigenvector, so
\begin{equation}
 \sum_xd_xg_x=\langle D^{1/2}\one,f\rangle=0.
 \label{eq:bounded-return-mean-zero}
\end{equation}
Moreover,
\begin{equation}
 \sum_xd_xg_x^2=\|f\|_2^2=F,
 \qquad
 \sum_{\{x,y\}\in E}(g_x-g_y)^2
 =f^\top Lf\le tF.
 \label{eq:bounded-return-norm-energy}
\end{equation}

By \eqref{eq:bounded-return-mean-zero} and $g_v>0$, some vertex $z$ has
$g_z\le0$.  Take a shortest path
$v=x_0,x_1,\ldots,x_{\ell}=z$ and set
$R:=\lfloor t^{-1/2}\rfloor\ge1$.

If $\ell\le R$, Cauchy--Schwarz along the path gives
\[
 \frac{g_v^2}{R}
 \le\frac{(g_v-g_z)^2}{\ell}
 \le\sum_{i=0}^{\ell-1}(g_{x_i}-g_{x_{i+1}})^2
 \le tF.
\]
Using \eqref{eq:bounded-return-projector-at-v} yields
$F\le d_vtR\le d_v\sqrt t$.

Suppose next that $\ell>R$.  If some $j\le R$ satisfies
$g_{x_j}\le g_v/2$, the same path argument on the prefix gives
\[
 \frac{g_v^2}{4R}\le tF,
\]
and hence $F\le4d_vtR\le4d_v\sqrt t$.  Otherwise
$g_{x_j}>g_v/2$ for every $0\le j\le R$.  Because the graph is unweighted
and has no isolated vertices, $d_{x_j}\ge1$, and therefore
\[
 F=\sum_xd_xg_x^2
 \ge\sum_{j=0}^{R}g_{x_j}^2
 >\frac{R+1}{4}g_v^2
 =\frac{R+1}{4}\frac{F^2}{d_v}.
\]
Thus $F<4d_v/(R+1)<4d_v\sqrt t$.  These cases prove
\eqref{eq:bounded-return-spectral-envelope}.
\end{proof}

\begin{remark}[Why the degree is absolute]
\label{rem:bounded-return-absolute-degree}
The factor $d_v$ enters through
$g_v=F_v(t)/\sqrt{d_v}$.  The proof also uses $d_x\ge1$ along a path.  Thus
\cref{lem:bounded-return-spectral-measure} is stated for ordinary degrees on
simple unweighted graphs.  A weighted extension would require explicit lower
control on path conductances and vertex measures; no such extension is
claimed.
\end{remark}

\section{The discounted diagonal and semantic support}
\label{sec:bounded-return-diagonal}

\begin{theorem}[Universal seed-return bound; \statusproved]
\label{thm:bounded-return-diagonal}
For every graph and seed in the scope of \cref{sec:bounded-return-scope},
\begin{equation}
 \boxed{
 \mathcal A_v
 \le
 \frac{d_v}{\vol(G)\gamma}
 +\frac{1}{1+c}
 +\frac{2\pi d_v}{\sqrt{c\gamma}}.}
 \label{eq:bounded-return-diagonal-bound}
\end{equation}
\end{theorem}

\begin{proof}
The matrices $H$ and $\gamma I+cL$ are diagonally similar, so their
inverse diagonal entries agree.  Hence
\begin{equation}
 \mathcal A_v
 =\sum_{k=1}^n\frac{\varphi_k(v)^2}{\gamma+c\lambda_k}
 =\frac{d_v}{\vol(G)\gamma}
  +\int_{(0,2]}\frac{1}{\gamma+ct}\,dF_v(t).
 \label{eq:bounded-return-spectral-representation}
\end{equation}
The first term is exactly the finite-graph stationary mode, not an error term.
Stieltjes integration by parts gives
\begin{equation}
 \int_{(0,2]}\frac{1}{\gamma+ct}\,dF_v(t)
 =\frac{F_v(2)}{\gamma+2c}
  +\int_0^2\frac{cF_v(t)}{(\gamma+ct)^2}\,dt.
 \label{eq:bounded-return-stieltjes}
\end{equation}
Here $F_v(0)=0$ by connectedness, $F_v(2)\le1$, and
$\gamma+2c=1+c$.  Applying
\cref{lem:bounded-return-spectral-measure} and enlarging the integral domain,
\[
 \int_0^2\frac{cF_v(t)}{(\gamma+ct)^2}\,dt
 \le4d_v\int_0^\infty\frac{c\sqrt t}{(\gamma+ct)^2}\,dt.
\]
With $t=\gamma x/c$,
\begin{equation}
 \int_0^\infty\frac{c\sqrt t}{(\gamma+ct)^2}\,dt
 =\frac{1}{\sqrt{c\gamma}}
   \int_0^\infty\frac{\sqrt x}{(1+x)^2}\,dx
 =\frac{\pi}{2\sqrt{c\gamma}}.
 \label{eq:bounded-return-beta-integral}
\end{equation}
Substitution into \eqref{eq:bounded-return-spectral-representation} proves
\eqref{eq:bounded-return-diagonal-bound}.  Eigenvalues near $-1$ for the
nonlazy random walk correspond to $\lambda_k$ near $2$ and are covered by
the bounded endpoint term; no aperiodicity assumption is used.
\end{proof}

\begin{corollary}[Bounded degree forbids simultaneous star scaling;
\statusproved]
\label{cor:bounded-return-support}
Suppose
\begin{equation}
 \vol(S_{\epsppr})\ge\frac{\sigma}{\epsppr}
 \label{eq:bounded-return-support-lower}
\end{equation}
for some $\sigma>0$.  Then
\begin{equation}
 \boxed{
 \alpha\mathcal A_v
 \le
 \frac{d_v(1+\alpha)\epsppr}{2\sigma}
 +\frac{\alpha(1+\alpha)}{2}
 +\frac{2\pi d_v(1+\alpha)\sqrt\alpha}
        {\sqrt{2(1-\alpha)}}.}
 \label{eq:bounded-return-support-bound}
\end{equation}
Consequently, for every joint sequence with
$\alpha\to0$, $\epsppr\to0$, $d_v\le D$, and
$\vol(S_{\epsppr})=\Theta(1/\epsppr)$,
\begin{equation}
 \alpha\frac{\pi_v}{\gamma_{\alpha}}\longrightarrow0.
 \label{eq:bounded-return-impossibility}
\end{equation}
In particular, such a family cannot have
$\pi_v/\gamma_{\alpha}=\Theta(1/\alpha)$.
\end{corollary}

\begin{proof}
The graph volume is at least the level-set volume, so
\eqref{eq:bounded-return-support-lower} implies
$1/\vol(G)\le\epsppr/\sigma$.  Substitute this into
\eqref{eq:bounded-return-diagonal-bound}, use
\[
 \gamma=\frac{2\alpha}{1+\alpha},
 \qquad
 c\gamma=\frac{2\alpha(1-\alpha)}{(1+\alpha)^2},
 \qquad
 \frac{1}{1+c}=\frac{1+\alpha}{2},
\]
and multiply by $\alpha$.  The resulting expression is
\eqref{eq:bounded-return-support-bound}, whose right side tends to zero under
the stated joint limit.
\end{proof}

\begin{corollary}[Necessary ordinary seed degree; \statusproved]
\label{cor:bounded-return-degree}
Under \eqref{eq:bounded-return-support-lower}, if
$\mathcal A_v\ge K/\alpha$, then, whenever the numerator is positive,
\begin{equation}
 \boxed{
 d_v\ge
 \frac{K-\alpha(1+\alpha)/2}
 {(1+\alpha)\epsppr/(2\sigma)
  +2\pi(1+\alpha)\sqrt\alpha/\sqrt{2(1-\alpha)}}.}
 \label{eq:bounded-return-degree-exact}
\end{equation}
Thus
\begin{equation}
 d_v=\Omega\!\left(\frac{1}{\epsppr+\sqrt\alpha}\right).
 \label{eq:bounded-return-degree-order}
\end{equation}
If $\alpha=O(\epsppr^2)$, this strengthens to
$d_v=\Omega(1/\epsppr)$.  If one also requires semantic nontriviality at the
seed, $\epsppr d_v\le\pi_v/2$, then $d_v=O(1/\epsppr)$, so in that parameter
regime $d_v=\Theta(1/\epsppr)$.
\end{corollary}

\begin{proof}
Rearrange \eqref{eq:bounded-return-support-bound} using
$\alpha\mathcal A_v\ge K$.  For fixed positive $K$ and sufficiently small
$\alpha$, the numerator is bounded below by a positive constant, giving
\eqref{eq:bounded-return-degree-order}.  When
$\sqrt\alpha=O(\epsppr)$, its denominator is $O(\epsppr)$.  Finally,
$\epsppr d_v\le\pi_v/2\le1/2$ gives the stated upper delimiter.
\end{proof}

\begin{remark}[The regime qualification is essential]
\label{rem:bounded-return-regime}
When $\epsppr=o(\sqrt\alpha)$,
\cref{cor:bounded-return-degree} yields only
$d_v=\Omega(1/\sqrt\alpha)$.  Therefore the unqualified chain
``$\Theta(1/\alpha)$ amplification, hence
$\Theta(1/\epsppr)$ seed degree'' is not established.  The latter conclusion
is correct from this argument only after imposing
$\alpha=O(\epsppr^2)$ and the seed nontriviality upper delimiter.
\end{remark}

\section{Closed-form family screen}
\label{sec:bounded-return-families}

The following formulas concern infinite rooted graphs.  Truncating beyond the
displayed $S_{\epsppr}$ scale gives finite graph families with the same
asymptotics and avoids any use of an infinite graph in the theorem.

\subsection{Endpoint path}

Put
\begin{equation}
 \lambda:=\frac{1-\sqrt\alpha}{1+\sqrt\alpha}.
 \label{eq:bounded-return-path-decay}
\end{equation}
On the endpoint-seeded half-line, direct solution of
$(D-cA)q=\gamma e_v$, $q=D^{-1}\pi$, gives
\begin{equation}
 q_j=\sqrt\alpha\,\lambda^j,
 \qquad
 \mathcal A_v=\frac{1}{\sqrt{1-c^2}}
 =\frac{1+\alpha}{2\sqrt\alpha}.
 \label{eq:bounded-return-path-formula}
\end{equation}
For fixed $\theta\in(0,1)$, take
$\epsppr=\theta\sqrt\alpha$.  If
$J=\max\{j:\lambda^j>\theta\}$, then
\[
 \vol(S_{\epsppr})=1+2J
 \sim\frac{\log(1/\theta)}{\sqrt\alpha}
 =\Theta(1/\epsppr).
\]
Thus the path saturates semantic support volume while amplifying only
$\Theta(\alpha^{-1/2})$.

\subsection{Caterpillar}

Consider a half-line backbone with one pendant leaf at every backbone vertex,
seeded at the endpoint.  Eliminating the leaves gives backbone decay
$q_j=q_0r^j$, where
\begin{equation}
 r=\frac{3-c^2-\sqrt{(3-c^2)^2-4c^2}}{2c},
 \qquad
 \mathcal A_v=\frac{2}{2-c^2-cr}.
 \label{eq:bounded-return-caterpillar-formula}
\end{equation}
As $\gamma\to0$,
\begin{equation}
 r=1-2\sqrt\gamma+O(\gamma),
 \qquad
 \mathcal A_v\sim\gamma^{-1/2}\sim(2\alpha)^{-1/2},
 \qquad
 q_0\sim\frac{\sqrt\gamma}{2}.
 \label{eq:bounded-return-caterpillar-asymptotic}
\end{equation}
A threshold equal to a fixed fraction of $q_0$ again gives
$\vol(S_{\epsppr})=\Theta(1/\epsppr)$.

\subsection{Hub ladder with bounded seed degree}

Let every backbone vertex carry one hub of degree $D$, with the remaining
$D-1$ hub neighbors pendant leaves.  The seed is the backbone endpoint and
has ordinary degree two, independent of $D$.  Eliminating the hub and its
leaves gives
\begin{equation}
 t:=\frac{c}{D-(D-1)c^2},
 \qquad
 r+r^{-1}=\frac{3-ct}{c},
 \qquad
 \mathcal A_v=\frac{2}{2-c(t+r)},
 \label{eq:bounded-return-hub-ladder-formula}
\end{equation}
where $r<1$ is the smaller root.  For fixed $D$,
\begin{equation}
 \mathcal A_v\sim\sqrt{\frac{2}{(D+1)\gamma}}
 \sim\frac{1}{\sqrt{(D+1)\alpha}},
 \qquad
 q_0\sim\sqrt{\frac{\gamma}{2(D+1)}}.
 \label{eq:bounded-return-hub-ladder-asymptotic}
\end{equation}
The active backbone length at a fixed-fraction threshold is
$\Theta((D\gamma)^{-1/2})$, while each module has volume $\Theta(D)$.
Therefore its semantic support still satisfies
\[
 \vol(S_{\epsppr})=\Theta(\sqrt{D/\gamma})
 =\Theta(1/\epsppr),
\]
but the hidden hub suppresses rather than enhances seed return.

\subsection{Trees and regular gadgets}

On the infinite $d$-regular tree,
\begin{equation}
 \mathcal A_v
 =\frac{2(d-1)}
 {d-2+\sqrt{d^2-4(d-1)c^2}}.
 \label{eq:bounded-return-regular-tree}
\end{equation}
For $d\ge3$, this tends to $(d-1)/(d-2)$; for $d=2$, it reduces to the path
formula.  On the infinite binary tree whose root has degree two and whose
other internal vertices have degree three,
\begin{equation}
 \mathcal A_v=\frac{4}{1+\sqrt{9-8c^2}}\longrightarrow2.
 \label{eq:bounded-return-binary-tree}
\end{equation}
Thus transient regular trees have bounded amplification.  Bounded-width
periodic ladders retain the one-dimensional $\Theta(\alpha^{-1/2})$ scale.
For a finite regular gadget with normalized spectral gap
$\lambda_2\ge\delta>0$, the exact spectral representation directly yields
\begin{equation}
 \mathcal A_v
 \le\frac{d_v}{\vol(G)\gamma}
 +\frac{1-d_v/\vol(G)}{\gamma+c\delta},
 \label{eq:bounded-return-expander}
\end{equation}
so only its finite-graph stationary term can scale as $1/\alpha$.

\subsection{Deterministic numerical screen}

The companion \texttt{verify.py} solves $(D-cA)q=\gamma e_v$ to relative
residual at most $10^{-11}$.  At every cell it selects a coordinate
breakpoint maximizing
$\epsppr\vol(S_{\epsppr})$, so the support column tests the most favorable
semantic threshold rather than a fixed arbitrary one.  Representative output
is shown in \cref{tab:bounded-return-screen}; it is \statusempirical{} and is
not used in the proof.

\begin{table}[t]
\centering
\small
\begin{tabular}{lcccc}
\toprule
Family ($d_v$) & $\alpha$ & $\alpha\mathcal A_v$
& $\sqrt\alpha\mathcal A_v$
& $\max_{\epsppr}\epsppr\vol(S_{\epsppr})$ \\
\midrule
Path (1) & $2^{-12}$ & $0.00781$ & $0.500$ & $0.374$ \\
Caterpillar (2) & $2^{-12}$ & $0.0109$ & $0.700$ & $0.376$ \\
Hub ladder, $D=32$ (2) & $2^{-12}$ & $0.00253$ & $0.162$ & $0.401$ \\
Finite binary tree (2) & $2^{-12}$ & $0.00061$ & $0.0388$ & $0.988$ \\
Three-regular cycle--antipode (3) & $2^{-12}$ & $0.00489$ & $0.313$ & $0.375$ \\
\bottomrule
\end{tabular}
\caption{Every screened bounded-degree family has vanishing
$\alpha\mathcal A_v$.  The one-dimensional families simultaneously keep the
semantic support-saturation score bounded away from zero.}
\label{tab:bounded-return-screen}
\end{table}

For comparison, on the centre-seeded star $K_{1,m}$,
\begin{equation}
 \mathcal A_v=\frac{(1+\alpha)^2}{4\alpha}.
 \label{eq:bounded-return-star}
\end{equation}
Taking $m=\Theta(1/\epsppr)$ makes every star vertex lie above a constant
multiple of $\epsppr$, so
$\vol(S_{\epsppr})=\Theta(1/\epsppr)$; its seed degree is exactly the large
quantity excluded by \cref{cor:bounded-return-support}.

\section{What is and is not closed}
\label{sec:bounded-return-scope-limit}

\paragraph{Closed.}
No family with bounded ordinary seed degree can simultaneously have
$\vol(S_{\epsppr})=\Theta(1/\epsppr)$ and
$\pi_v/\gamma_{\alpha}=\Theta(1/\alpha)$ in the joint local limit.  Hence a
signed-relaxation seed-operation lower bound whose only large instance
parameter is $\pi_v/\gamma_{\alpha}$ cannot reproduce the star's
$\Theta(\alpha^{-1/2})$ number of seed operations on such a family.

\paragraph{Not closed.}
An upper bound on one Green diagonal is not an algorithmic upper bound.  A
different graph may force work at many vertices, or couple several modes so
that no single seed-travel certificate captures the total charged work.
Therefore this note does not rule out
$\Omega(1/(\sqrt\alpha\,\epsppr))$ work for every member of a signed one-hop
class on a bounded-degree family.  It also says nothing about stronger block,
elimination, or response primitives.  The next valid lower-bound target is a
distributed obstruction, not another bounded-degree self-return search.

\end{document}
